\documentclass[11pt]{article}
\usepackage[T1]{fontenc}
\usepackage{amsmath,amssymb,amsthm}
\usepackage[margin=1in]{geometry}
\usepackage{booktabs}
\usepackage[hidelinks]{hyperref}
\hypersetup{
  pdftitle={A Conjectural Unary-Theta Formula for the 6n+1-Section of 18-Colored Generalized Frobenius Partitions},
  pdfauthor={Xiaojian Zeng},
  pdfsubject={18-colored generalized Frobenius partitions and a conjectural modulo-16 identity},
  pdfkeywords={generalized Frobenius partitions, congruences, q-series, exact computation}
}

\newtheorem{proposition}{Proposition}

\newcommand{\cphi}{c\phi}
\newcommand{\CT}{\operatorname{CT}}

\title{A Conjectural Unary-Theta Formula for the $6n+1$-Section of\\
18-Colored Generalized Frobenius Partitions}
\author{Xiaojian Zeng\\
PBC School of Finance, Tsinghua University, 100190\\
\texttt{xiaojianzeng@foxmail.com}}
\date{}

\begin{document}
\maketitle

\begin{abstract}
Das, Nath, and Sarma recently conjectured two congruences modulo $16$ for
the $18$-colored generalized Frobenius partition function.  We propose the
stronger, explicitly unproved, coefficientwise congruence
\[
 \sum_{m\geq0}\cphi_{18}(6m+1)q^m
 \equiv 4\frac{f_2f_3^2}{f_1f_6}
 =4\sum_{k\in\mathbb Z}q^{k(3k-1)/2}\quad(\mathrm{mod}\ 16),
 \qquad f_d=(q^d;q^d)_\infty.
\]
We prove the equality between the two expressions on the right and show that
the proposed congruence would imply both source congruences.  Exact
computations from two different generating descriptions support it for
$0\leq m\leq504$; the range $0\leq m\leq99$ was used for discovery and
$100\leq m\leq504$ was held out.  A termwise dissection audit gives several
formal parity cancellations but leaves other grouped cancellations unproved.
Thus neither the displayed congruence nor the original problem is resolved.
\end{abstract}

\section{Statement and context}

A generalized Frobenius partition is an equal-length two-row array of
nonnegative integers whose weight is the row length plus the sum of its
entries.  Andrews introduced the colored version, in which the entries are
taken from fixed ordered copies of the nonnegative integers
\cite{Andrews1984GFP}.  We write \(\cphi_k(N)\) for the number of such
partitions of \(N\) with \(k\) colors.  Equivalently, the generating function
needed below is
\begin{equation}\label{eq:lattice}
 \sum_{N\geq0}\cphi_{18}(N)q^N
 =\frac{1}{(q;q)_\infty^{18}}
   \sum_{(x_1,\ldots,x_{17})\in\mathbb Z^{17}}
   q^{Q(x_1,\ldots,x_{17})},
\end{equation}
where
\[
 Q(x)=\sum_{i=1}^{17}x_i^2+
       \sum_{1\leq i<j\leq17}x_ix_j.
\]
Constant-term descriptions and integral expressions for colored generalized
Frobenius generating functions were developed by Cui, Gu, and Tang
\cite{CuiGuTang2024,CuiGuTang2025}.  These descriptions are especially useful
when exact coefficients, rather than floating-point approximations, are
required.

Das, Nath, and Sarma proved congruences for \(\cphi_{18}\) modulo $8$,
$81$, and $2187$, and then stated
\begin{equation}\label{eq:source-conj}
 \cphi_{18}(30n+19)\equiv0\quad(\mathrm{mod}\ 16),\qquad
 \cphi_{18}(30n+25)\equiv0\quad(\mathrm{mod}\ 16)
 \quad(n\geq0)
\end{equation}
as a numerical conjecture whose proof was left for future investigation
\cite[Conjecture 5.1]{DasNathSarma2026}.  Their nearby modulus-$8$ theorem is
on the class $6n+5$; it does not apply directly to either class in
\eqref{eq:source-conj}, since both $30n+19$ and $30n+25$ are $1$
modulo $6$.

Put
\[
 H(q)=\sum_{m\geq0}\cphi_{18}(6m+1)q^m.
\]
All congruences of power series below are coefficientwise.  Our proposed
refinement is the following.

\medskip
\noindent\textbf{Conjecture.}
In \((\mathbb Z/16\mathbb Z)[[q]]\),
\begin{equation}\label{eq:main-conj}
 H(q)\equiv4\frac{f_2f_3^2}{f_1f_6}\quad(\mathrm{mod}\ 16).
\end{equation}
The quotient is well defined as a formal power series because every $f_d$
has constant coefficient one.  The conjecture is stronger than
\eqref{eq:source-conj}: it predicts every coefficient in the entire
$6n+1$ section modulo $16$, rather than only two residue classes within
that section.  A bounded search of the cited sources and exact formula
searches through September 7, 2026, located no source stating
\eqref{eq:main-conj}; this is a description of the search, not a priority
claim.

\section{A proved identity and the implication for the source conjecture}

The right side of \eqref{eq:main-conj} has a sparse description that explains
the two progressions in \eqref{eq:source-conj}.

\begin{proposition}\label{prop:theta}
In \(\mathbb Z[[q]]\),
\begin{equation}\label{eq:theta}
 \frac{f_2f_3^2}{f_1f_6}
 =\sum_{k\in\mathbb Z}q^{k(3k-1)/2}.
\end{equation}
Consequently, if \eqref{eq:main-conj} holds, then both congruences in
\eqref{eq:source-conj} hold for every \(n\geq0\).
\end{proposition}

\begin{proof}
The Jacobi triple product in the form recorded, for example, in
\cite[Theorem 2.8]{Andrews1998Partitions}, with base \(q^3\) and parameter
\(q\), gives
\[
 \sum_{k\in\mathbb Z}q^{k(3k-1)/2}
 =f_3(-q;q^3)_\infty(-q^2;q^3)_\infty.
\]
Using $1+x=(1-x^2)/(1-x)$ in the two nonzero residue classes modulo $3$
turns the product on the right into
\[
 f_3\prod_{r\geq1,\,3\nmid r}
 \frac{1-q^{2r}}{1-q^r}
 =f_3\frac{f_2/f_6}{f_1/f_3}
 =\frac{f_2f_3^2}{f_1f_6},
\]
which proves \eqref{eq:theta} formally.

For $k$ modulo $5$, the exponent $k(3k-1)/2$ takes only the residues
$0,1,2$.  Thus the coefficients of $q^{5n+3}$ and $q^{5n+4}$ on the
right of \eqref{eq:theta} vanish.  Under \eqref{eq:main-conj} this yields
\[
 \cphi_{18}(6(5n+3)+1)\equiv
 \cphi_{18}(6(5n+4)+1)\equiv0\quad(\mathrm{mod}\ 16),
\]
and the two arguments are $30n+19$ and $30n+25$, respectively.
\end{proof}

\section{Exact finite evidence}

The finite test covered all coefficients of the full generating series
through $q^{3025}$, hence all coefficients of $H(q)$ through $m=504$.
The first computation used the direct Laurent constant term
\begin{equation}\label{eq:ct}
 \CT_z\,(1+z^{-1})^{18}
 \prod_{w\geq1}(1+zq^w)^{18}(1+z^{-1}q^w)^{18}.
\end{equation}
For degrees at most $3025$, factors with $w>3025$ can be omitted.  The
largest necessary row length is $233$: lengths $233$ and $234$ have
minimum possible weights $3017$ and $3042$, respectively.  This makes the
Laurent truncation at absolute $z$-degree $233$ exhaustive in the stated
range.  Arithmetic was performed modulo $256$ and then reduced modulo
$16$, which is exact because the reduction map is a ring homomorphism.

The second computation independently expanded formula (3.15) of
Das, Nath, and Sarma \cite{DasNathSarma2026} modulo $16$,
retaining fourteen summands. It used Euler's pentagonal theorem,
the partition recurrence for reciprocal eta products, and enumeration of the
Borwein cubic theta series.  The two resulting $3026$-term residue vectors
agreed.  On the right of \eqref{eq:main-conj}, one calculation expanded the
eta quotient, while another directly enumerated the generalized pentagonal
exponents.  All four section vectors agreed for $0\leq m\leq504$.

The support pattern was recognized on $0\leq m\leq99$, after which
\eqref{eq:main-conj} was fixed and tested on the held-out range
$100\leq m\leq504$.  In particular, the finite computation verifies both
congruences in \eqref{eq:source-conj} only for $0\leq n\leq100$.  It does
not establish any case beyond that bound and supplies no finite-to-infinite
reduction.  This separation is the same basic scope issue that arises in
exact finite-certificate work on restricted models
\cite{ZengEtAl2026Finite}: exhaustive exactness inside a stated finite model
does not settle its ambient open problem.

\section{Structural attack and remaining gap}

A termwise audit of formula (3.15) \cite{DasNathSarma2026} first deletes
the four summands with multipliers divisible by $16$. Number the fourteen
remaining summands in their original displayed order. The
fourth, fifth, sixth, tenth, and twelfth summands carry a factor $8$ modulo
$16$.  After that factor is removed and the remaining eta quotients are
reduced modulo $2$, their residual products are, respectively,
\[
 \frac{f_{48}}{f_2f_{16}},\quad
 \frac{f_{48}}{f_2f_{16}},\quad
 \frac{f_6f_{48}}{f_8f_{16}},\quad
 \frac{f_6f_{96}}{f_8f_{16}},\quad
 \frac{f_6f_{96}}{f_8f_{16}}.
\]
Their shifts are $q^2,q^2,q^2,q^4,q^4$.  Every displayed product has only
even powers of $q$, so these five terms make no contribution to any odd
coefficient, formally and to all orders.  The same modulo-$2$ reductions
also prove the whole-series cancellations of the pairs consisting of terms
$(4,5)$ and $(10,12)$ modulo $16$.

In the tested $6m+1$ section, the remaining summands split into the groups
\[
 (1+2),\qquad(3+7),\qquad(8+9),\qquad(11+13+14).
\]
The last three groups vanished for every tested $m$, while the first group
agreed with $4\sum_kq^{k(3k-1)/2}$ and had zero fifth-sections in residues
$3$ and $4$.  These grouped statements are finite observations through
$m=504$, not formal identities.  They are the unresolved kernel of the
proposed congruence.

A direct continuation is to apply the relevant $2$-dissections to formula
(3.15) modulo $16$, extract odd powers, and then apply $3$-dissections to
the resulting indices congruent to $0$ modulo $3$.  After the five formal
parity eliminations above, one must prove that the four surviving groups
combine to $4f_2f_3^2/(f_1f_6)$.  A modular-form alternative would require
an explicit level, character, integral coefficient ring, cusp-holomorphy or
denominator-clearing argument, and a justified Sturm-type bound; none is
provided here.

The computational conjecture is immediately falsifiable: extend either exact
left-side generator beyond degree $3025$, and compare
\(\cphi_{18}(6m+1)\bmod16\) with $4$ when $m=k(3k-1)/2$ for some
$k\in\mathbb Z$, and with $0$ otherwise.  Any mismatch disproves
\eqref{eq:main-conj}.  The supplied scripts replay the present finite test via
\begin{verbatim}
python3 evidence/computation/constant_term_mod16.py
python3 evidence/computation/source_formula_mod16.py
python3 evidence/computation/h_identity_mod16.py
python3 evidence/computation/projection_term_audit.py
\end{verbatim}

\section*{Computational and AI assistance}

All reported coefficient calculations used exact integer or modular
arithmetic; no floating-point comparison was used.  Python and NumPy were
used to implement the two left-side generators, the two right-side
generators, and the summand audit.  An OpenAI language model assisted with
formula exploration, draft code, adversarial checks, literature screening,
and initial manuscript preparation.  Its output is not mathematical
evidence.  In particular, the language model did not prove
\eqref{eq:main-conj}; the universal identity and both source congruences
remain unresolved.

\bibliographystyle{abbrv}
\bibliography{references}

\end{document}
